\input{preamble.tex}

\newcommand{\DocTitle}{Pointer Authentication on arm64e: tightening function pointers}
\newcommand{\DocDescription}{one valid pointer in the wrong slot}
\newcommand{\DocTags}{writeup, cybersecurity, low-level, arm64e}

\begin{document}

\webimgw{/img/PAC-lldb.png}{LLDB showing a signed function pointer, its PAC field, and the actual target address}{}{1000}

Imagine this, the compiler loads a function pointer from memory and branches to whatever address it finds:

\codelang{asm}
\begin{verbatim}
; plain arm64, -O0
ldr  x8, [sp, #16]   ; load the stored function pointer
blr  x8              ; branch to it; no integrity check
\end{verbatim}

If an attacker can write that memory slot, they choose the branch target.

We demonstrated the redirect in \texttt{m1\uscore{}fnptr.c}: overwriting \texttt{o.op} made the program call \texttt{unexpected\uscore{}handler} instead of \texttt{add\uscore{}one}. On arm64e the same call site is different:

\codelang{asm}
\begin{verbatim}
; arm64e, -O0
ldr    x8, [sp, #16]
blraaz x8            ; authenticated indirect branch (IA key, zero modifier)
\end{verbatim}

The CPU will now reject a pointer whose signature does not verify.

Pointer Authentication (PAC) turns a pointer into a value that the CPU can verify before use. On AArch64 it is an architecture extension (PAuth, mandatory from Armv8.3-A); on Apple platforms it is deployed as the \textbf{arm64e} ABI. PAC is widely described as a mitigation for control-flow hijacking, but the strength of that mitigation depends entirely on \textit{which authentication context} a pointer is bound to. This is what I mean:

The default arm64e context for C function pointers is intentionally broad:

\begin{itemize}
\item one key (for different instruction / data pointers)
\item discriminator zero (acts like associated data in encryption, but this is purely for hashing)
\item no address diversity
\end{itemize}

This begs the question, how can we ensure that a function pointer is only usable in its intended authentication context, rather than relying exclusively on the default pointer authentication scheme?

The goal of this my exploration into PAC is to be able to understand and find an answer to this question. Particularly, this involves how we might tighten scopes with the \texttt{\uscore{}\uscore{}ptrauth} type qualifier on MacOS.

\section{How Pointer Authentication Works}

\textbf{A PAC is a keyed hash, not encryption.} \texttt{sign(pointer, modifier, key)} computes a MAC (Message Authentication Code) over the pointer and a 64-bit modifier, then stores the truncated result in the pointer's spare high bits. \texttt{authenticate} recomputes it, compares, and removes it.

At first glance, it might seem useless to compute a long MAC that just gets truncated because that would reduce the number of possible bit combinations an attacker must try, but since we store the PAC alongside the 64 bits of the pointer itself, we need to truncate it.

Computing the longer MAC is still very important because it allows for greater entropy in the result that still gets carried on to the truncated PAC.

\subsection{Keys}

Arm PAuth provides five 128-bit keys, stored in system registers that only privileged code can read or write:

\begin{tabular}{|l|l|}
\hline
Key & Use \\
\hline
\texttt{IA} (ASIA, 0) & instruction pointers \\
\hline
\texttt{IB} (ASIB, 1) & instruction pointers \\
\hline
\texttt{DA} (ASDA, 2) & data pointers \\
\hline
\texttt{DB} (ASDB, 3) & data pointers \\
\hline
\texttt{GA} (ASGA) & generic \texttt{pacga} signatures \\
\hline
\end{tabular}

Which key signs which pointer is an \textbf{ABI} decision which can vary across architectures.

One thing that confused me was that A and B keys do not have the same uses across operating systems.

Key usage is an ABI thing, not a hardware thing. In our generic AArch64 build with \texttt{-march=armv8.3-a -mbranch-protection=pac-ret}, Clang signed the return address with IA. On arm64e, it uses IB.

Clang's docs say arm64e mostly uses A keys for global stuff like function pointers and vtables, and B keys for local stuff like return addresses. A keys tend to be shared across processes, B keys tend not to be. That's an OS convention, not a hardware promise.

\subsection{Modifiers / discriminators}

The modifier is the extra 64-bit value mixed into the hash. Clang calls it a discriminator.

\subsection{Sign, authenticate, strip}

Just an FYI, throughout this writeup, I might put a * after things like PAC, AUT, or XPAC. When I do this, I am referring to a wild card. I know with pointers * means dereference and that can be confusing in this context, so I just want to get that settled before I confuse anyone haha. Alright, moving on:

\begin{itemize}
\item \texttt{PAC*} signs. Observed form: \texttt{pacda x8, x9} (sign the data pointer in \texttt{x8}, modifier in \texttt{x9}, DA key).
\item \texttt{AUT*} authenticates. Baseline Armv8.3 \texttt{AUT*} does \textbf{not} trap on failure; it leaves a corrupted pointer that faults when used.
\item \texttt{XPAC*} strips the code without authenticating.
\item \texttt{FEAT\uscore{}FPAC} (Armv8.6-A, "Faulting on PAC") makes a failed \texttt{AUT*} raise an exception immediately.
\end{itemize}

On baseline Armv8.3, \texttt{AUT*} does not raise an exception. On failure it returns a corrupted pointer that would only fault when that pointer is later used.

Clang therefore adds its own check and trap around it, which I found below at \texttt{-O2}:

\codelang{asm}
\begin{verbatim}
autda  x16, x8            ; authenticate (DA key), modifier in x8
mov    x17, x16
xpacd  x17                ; strip the PAC
cmp    x16, x17
b.eq   .Lauth_success
brk    #0xc472            ; trap if authentication failed
.Lauth_success:
ldr    w0, [x16]
\end{verbatim}

On success the authenticated value is already canonical, so stripping it is a no-op and the compare succeeds. On failure the corrupted high bits remain in \texttt{x16}, so the compare fails.

The \texttt{brk} is not part of \texttt{AUT*}. It is the explicit trap Clang inserted \textit{precisely because} \texttt{AUT*} does not trap. With \texttt{FEAT\uscore{}FPAC}, the hardware faults at \texttt{AUT*} itself and this check is unnecessary.

\subsection{Where the PAC lives (measured on the test host)}

I probably could have just looked this up for my host architecture, but I wanted to have a little fun:

I signed 1024 data pointers with different discriminators, XOR each signed value with its raw address to cancel the address bits (\texttt{uintptr\uscore{}t diff = signed \^{} raw}), then bitwise-OR all those \texttt{diff} values together (\texttt{mask |= diff}). The only bits that ever changed would be our PAC field:

\codelang{text}
\begin{verbatim}
union of changed bits: 0x007f800000000000
set bit positions: 54 53 52 51 50 49 48 47    -> 8-bit field
\end{verbatim}

So on this host the PAC occupies bits 47--54 (an 8-bit field).

You might be thinking now, 8 bits does not give many possible values. Could an attacker not try a blind forgery?

They could, but it's not a free brute force. A forged pointer has a 1/256 chance of carrying the right signature for that key, modifier, and target address. A wrong guess traps right away, either from FPAC or Clang's \texttt{brk}, so there's no retry loop in the same process and no useful feedback from failures.

The bigger point is substitution. Copying a validly signed pointer from one field into another only works when both fields use the same signing schema. Under arm64e's default \texttt{IA,0} schema, it works. That's the A case in \texttt{hardened\uscore{}ops.c}. Once each field has its own context, the same copy fails every time. That's the part I'm getting at later.

\section{Understanding arm64e and Its Default ABI}

Terms:

\begin{itemize}
\item \textbf{AArch64} the instruction set.
\item \textbf{PAuth / PAC} the optional AArch64 extension providing \texttt{PAC*}/\texttt{AUT*}.
\item \textbf{arm64e} Apple's ABI: a contract specifying which pointers are signed, with which key and modifier, how signed pointers are stored, and how the linker and dynamic loader cooperate.
\end{itemize}

\subsection{The default key map (arm64e)}

\begin{tabular}{|l|l|l|l|}
\hline
Pointer kind & Key & Modifier & Observed on the test host \\
\hline
Return address & IB & stack pointer on entry & \texttt{pacibsp} \ldots{} \texttt{retab} \\
\hline
C function pointer & IA & \textbf{always 0} & \texttt{paciza}, \texttt{@AUTH(ia,0)}, \texttt{braaz}/\texttt{blraaz} \\
\hline
C++ vtable pointer & DA & address diversity + vtable id & (runtime) \\
\hline
ObjC \texttt{isa}/\texttt{super} & DA & address diversity + \texttt{0x6AE1} / \texttt{0x25DA} & (runtime) \\
\hline
ObjC selector & DB & \texttt{0x57C2} for \texttt{SEL} ivars & (runtime) \\
\hline
Ordinary C data pointer & none & --- & raw \texttt{.quad}, no signature \\
\hline
\end{tabular}

Note: the vtable pointer and the function entries inside it don't use the same schema. The pointer uses \texttt{DA}, address diversity, and a discriminator from the mangled vtable identifier of the primary base class. The function entries use \texttt{IA}, address diversity, and a discriminator from the function's mangled name. For \texttt{isa} and \texttt{super}, both use address diversity, with \texttt{0x6AE1} for \texttt{isa} and \texttt{0x25DA} for \texttt{super}.

Measured, \texttt{m3\uscore{}defaults.c}:

\codelang{asm}
\begin{verbatim}
; return address
pacibsp
...
retab

; C function pointer in data and its call site
_fn_ptr:
    .quad _target@AUTH(ia,0)     ; signed at load: IA key, discriminator 0
...
adrp x8, _fn_ptr@PAGE
ldr  x1, [x8, _fn_ptr@PAGEOFF]
braaz x1                          ; authenticated branch, zero modifier

; ordinary data pointer (same file)
_data_ptr:
    .quad _data_global            ; raw address, no signature
\end{verbatim}

\subsection{The default C function-pointer context}

Key \texttt{IA}, discriminator \texttt{0}, no address diversity.

There is effectively \textbf{one} function-pointer context in the WHOLE process. Any correctly signed C function pointer can authenticate in any C function-pointer slot. The ABI gives us \textit{"is this a validly signed function pointer?"} but not \textit{"is it the function this field is supposed to hold?"}

This is not an oversight though, because C function pointers must remain interchangeable across static libraries as well as separately compiled objects like translation units and dyld shared caches. The ABI, therefore, cannot bake a per-field discriminator into every function pointer and break compatibility.

\subsection{arm64e is not "AArch64 with PAC"}

The same instructions can be spent by different ABIs in very different ways. A generic AArch64 (Linux) build with \texttt{-march=armv8.3-a -mbranch-protection=pac-ret} produced:

\codelang{asm}
\begin{verbatim}
caller:
    paciasp          ; return address signed with the IA key (A key by default)
    ...
    retaa

indirect:
    ldr  x2, [x0]
    br   x2          ; function pointers are NOT authenticated
\end{verbatim}

versus arm64e:

\codelang{asm}
\begin{verbatim}
caller:
    pacibsp          ; IB key
    ...
    retab

indirect:
    ldr  x1, [...]
    braaz x1         ; authenticated
\end{verbatim}

There are two main differences here.

First, the return-address key choice differs: generic pac-ret defaults to the A key whereas arm64e uses the B key.

Second, the base Linux ABI does not sign function pointers at all, while arm64e signs them automatically.

What Apple added on top of Arm's extension is therefore \textbf{policy and infrastructure}:

\begin{itemize}
\item an ABI specification (which pointer kinds are signed, with which keys and modifiers)
\item compiler, linker, and dyld integration
\item runtime support in objc4, libc++, and Swift, each with assigned discriminators
\item conventions for key usage (the A / B keys usages that I spoke of earlier in the article);
\end{itemize}

\section{Pointer Substitution and Authentication Contexts}

A pointer's authentication context has two inputs: the \textbf{key} and the \textbf{modifier}. The key is secret and the modifier is not.

To prevent substitution, the context must be different for each field we care about, so a pointer signed for one field cannot authenticate in another.

When the compiler knows both fields statically, it does the right thing on a typed assignment. It authenticates under the source schema and re-signs under the destination schema.

\section{Discriminators and Address Diversity}

Two independent ways to build a field-specific context:

\begin{itemize}
\item A \textbf{constant discriminator} (16-bit) distinguishes categories: \texttt{0x1234} for \texttt{open}, \texttt{0x5678} for \texttt{read}, \texttt{0x9ABC} for \texttt{close}.
\item \textbf{Address diversity} binds the pointer to its \textit{storage location} (the modifier includes the address at which the pointer is stored)
\end{itemize}

Blending the two means replacing the top 16 bits of the storage address with the constant. The compiler generates this, then signs it with the IA key:

\codelang{asm}
\begin{verbatim}
ldr  x1, [x0]                ; load pointer from the field
mov  x16, x0                 ; x0 = &struct (storage address)
movk x16, #4660, lsl #48     ; 4660 = 0x1234 -> top 16 bits
braa x1, x16                 ; authenticated branch, modifier in x16

; initialization for the same field:
mov  x16, x19                ; storage address
movk x16, #4660, lsl #48
pacia x8, x16                ; sign with IA key
\end{verbatim}

Because the context contains both the field's address and its discriminator, a value copied from \texttt{read} to \texttt{open} carries the wrong modifier and fails authentication.

Clang also rejects direct \texttt{memcpy} between differently-qualified pointer types at compile time (which is pretty cool if you ask me).

\section{Enforcing Stronger Authentication with \texttt{\uscore{}\uscore{}ptrauth}}

Okay FINALLY onto using \texttt{\uscore{}\uscore{}ptrauth} to give your code PAC where it might not be automatic.

\texttt{\uscore{}\uscore{}ptrauth} is a Clang type qualifier: \texttt{\uscore{}\uscore{}ptrauth(key, address\uscore{}diversity, discriminator)}. It goes after the * and before the field name. Include \texttt{<ptrauth.h>} to use the named keys:

\codelang{c}
\begin{verbatim}
#include <ptrauth.h>

#define DEVICE_OP(d) __ptrauth(ptrauth_key_function_pointer, 1, d)

struct DeviceOps {
    int (*DEVICE_OP(0x1234) open)(void);
    int (*DEVICE_OP(0x5678) read)(void);
    int (*DEVICE_OP(0x9ABC) close)(void);
};
\end{verbatim}

\texttt{ptrauth\uscore{}key\uscore{}function\uscore{}pointer} is the IA key in this toolchain. The \texttt{1} enables address diversity. The third argument is a 16-bit constant discriminator, unique per field.

\section{Building a Hardened Function-Pointer Table}

\texttt{hardened\uscore{}ops.c} builds both versions side by side and tests each substitution in a forked child so the failed authentication is contained:

\codelang{text}
\begin{verbatim}
A: raw substitution        accepted, returned 22
B: raw substitution        REJECTED (signal 10)
B: typed assignment        accepted, returned 22
\end{verbatim}

Version A is the ordinary table. version B applies the qualifiers above.

Reproduce with:

\codelang{sh}
\begin{verbatim}
clang -target arm64e-apple-macos11 -O1 hardened_ops.c -o hardened_ops
./hardened_ops
\end{verbatim}

\section{References}

\begin{itemize}
\item \href{https://clang.llvm.org/docs/PointerAuthentication.html}{Clang Pointer Authentication documentation}
\item \href{https://llvm.org/docs/PointerAuth.html}{LLVM PointerAuth backend documentation}
\item Clang \texttt{ptrauth.h} (key assignments and intrinsic contracts; shipped in the compiler resource directory).
\item \href{https://developer.arm.com/documentation/ddi0487/latest}{Arm Architecture Reference Manual for A-profile, Pointer Authentication}
\item \href{https://developer.apple.com/documentation/security/preparing-your-app-to-work-with-pointer-authentication}{Apple, "Preparing your app to work with pointer authentication"}
\end{itemize}

\end{document}
